\section{A4: Field operator and density operator}

For a complete orthonormal one-particle basis $\{\phi_j(\mathbf r)\}$, define
\[
\boxed{
\hat\Psi(\mathbf r)=\sum_j\phi_j(\mathbf r)\hat a_j,
\qquad
\hat\Psi^\dagger(\mathbf r)
=\sum_j\phi_j^*(\mathbf r)\hat a_j^\dagger.
}
\]
The field operator annihilates a boson at position $\mathbf r$; it is the
position-basis form of the mode annihilation operators.

Using orthonormality gives the inverse transformation:
\[
\boxed{
\hat a_i=\int d^3r\,\phi_i^*(\mathbf r)\hat\Psi(\mathbf r),
\qquad
\hat a_i^\dagger=\int d^3r\,\phi_i(\mathbf r)
\hat\Psi^\dagger(\mathbf r).
}
\]

The bosonic mode algebra and completeness imply
\[
\begin{aligned}
[\hat\Psi(\mathbf r),\hat\Psi^\dagger(\mathbf r')]
&=\sum_{ij}\phi_i(\mathbf r)\phi_j^*(\mathbf r')
[\hat a_i,\hat a_j^\dagger]\\
&=\sum_i\phi_i(\mathbf r)\phi_i^*(\mathbf r')\\
&=\delta^{(3)}(\mathbf r-\mathbf r').
\end{aligned}
\]
Also,
\[
[\hat\Psi(\mathbf r),\hat\Psi(\mathbf r')]=0,
\qquad
[\hat\Psi^\dagger(\mathbf r),\hat\Psi^\dagger(\mathbf r')]=0.
\]

For example,
\[
\hat\Psi(\mathbf r)|2,1\rangle
=\sqrt2\,\phi_0(\mathbf r)|1,1\rangle
+\phi_1(\mathbf r)|2,0\rangle.
\]
Removal at a position is a superposition of removal from every occupied
orbital that has amplitude there.

The local number-density operator is
\[
\boxed{
\hat n(\mathbf r)
=\hat\Psi^\dagger(\mathbf r)\hat\Psi(\mathbf r).
}
\]
For a Fock state,
\[
\langle\hat a_i^\dagger\hat a_j\rangle
=n_j\delta_{ij},
\]
so
\[
\boxed{
n(\mathbf r)=\langle\hat n(\mathbf r)\rangle
=\sum_jn_j|\phi_j(\mathbf r)|^2.
}
\]
In a completely condensed state $|N,0,\ldots\rangle$,
\[
n(\mathbf r)=N|\phi_0(\mathbf r)|^2.
\]

Integrating the local density recovers the total-number operator:
\[
\boxed{
\hat N=\int d^3r\,
\hat\Psi^\dagger(\mathbf r)\hat\Psi(\mathbf r)
=\sum_j\hat a_j^\dagger\hat a_j.
}
\]
Hence $[\hat\Psi]=L^{-3/2}$ in three dimensions.

The operator $\hat\Psi(\mathbf r)$ must be distinguished from the later
mean-field condensate wavefunction $\psi(\mathbf r)$. In a symmetry-breaking
description one eventually writes
$\psi(\mathbf r)=\langle\hat\Psi(\mathbf r)\rangle$, but this is an
additional approximation and interpretation.
